\centerline{\bf \Huge }


$\bigl[\textbf{КМО 2026, юниоры, вторая задача.}\bigr]$
В трапеции $ABCD$ диагонали $AC$ и $BD$ перпендикулярны. На основании $AD$ отмечена точка К такая, что $AK+KC = AD+ BC$. Докажите, что точка $K$ равноудалена от точек $A$ и $C$.


Пусть $O$ - точка пересечения диагоналей, и мы выбрали комплексную плоскость так, что ось $O x$ направлена вдоль $A C$, а ось $O y$ - вдоль $B D$. Тогда можно записать

$$
A=-a, \quad C=c, \quad B=i b, \quad D=-i d,
$$


где $a, b, c, d>0$.

То есть точки $A, C$ лежат на вещественной оси, а $B, D$ --- на мнимой.

$A D \| B C$ в комплексной записи означает, что число $\dfrac{a-i d}{c-i b}$ чисто вещественно. Значит, $Im((a-i d)(c+i b))=0$, то есть $ab - cd = 0$, то есть $\dfrac{d}{a} = \dfrac{b}{c} = L$ (на самом то деле это можно было понять просто из подобия, но наша цель как можно больше попрактиковаться с комплексными числами).

Параметризуем точку $K\in AD$. В комплексных числах это записывается следующим образом

$$
K=A+t(D-A)=-a+t(a-i d), \quad 0 \leq t \leq 1 .
$$


То есть

$$
K=a(t-1)-i d t
$$

Теперь выразим все нужные нам длины:

$$A K=t|D-A|=t \sqrt{a^2+d^2}= ta\sqrt{1+\dfrac{d^2}{a^2}}=t a \sqrt{1+L^2}$$.

\begin{flushright}
    \vspace{-20mm}
    (1)
\end{flushright}

Напомним, что $L$ этот тот самый коэффициент подобия $\dfrac{a}{c}=\dfrac{b}{d} = L$.

Далее,

$$
A D=\sqrt{a^2+d^2}=a \sqrt{1+L^2}, \quad B C=\sqrt{c^2+b^2}=c \sqrt{1+L^2},
$$

Подставим это всё в равенство $AK+KC = AD+BC$ и получим

$$
ta\sqrt{1+L^2} + KC = (a+c)\sqrt{1+L^2}
$$

отсюда 

$$
KC = (a+c-at)\sqrt{1+L^2}
$$. 

\begin{flushright}
    \vspace{-20mm}
    (2)
\end{flushright}
Но с другой стороны $K C^2=|K-C|^2=(a+c-a t)^2+d^2 t^2=(a+c-a t)^2+a^2 L^2 t^2$, так как $K-C = at - a - c - idt$. Приравняем квадрат первого выражения со вторым и получим:

$$
(a+c-a t)^2+a^2 L^2 t^2=(a+c-a t)^2\left(1+L^2\right)
$$

$$
a^2 L^2 t^2=L^2(a+c-a t)^2
$$

$$
a t=a+c-a t
$$

$$
at= \dfrac{a+c}{2}
$$

Подставив это значение $at$ в формулы $(1)$ и $(2)$ получаем, что $AK=KC = \dfrac{a+c}{2}\sqrt{1+L^2}$